% !TeX program = pdflatex
\documentclass[11pt]{amsart}
\usepackage[letterpaper,margin=1in]{geometry}
\usepackage{amssymb,mathtools,microtype,enumitem,booktabs}
\usepackage[square,numbers,sort&compress]{natbib}
\usepackage{xurl}
\usepackage[colorlinks=true,linkcolor=blue,citecolor=blue,urlcolor=blue]{hyperref}
\hypersetup{pdftitle={A general shape and fluctuation theorem for inward-excited walks},
pdfauthor={Ahmed Bou-Rabee, Tuan-Minh Nguyen, Yuval Peres, Arvind Singh}}
\theoremstyle{plain}
\newtheorem{theorem}{Theorem}[section]
\newtheorem{lemma}[theorem]{Lemma}
\newtheorem{proposition}[theorem]{Proposition}
\newtheorem{corollary}[theorem]{Corollary}
\theoremstyle{definition}
\newtheorem{remark}[theorem]{Remark}
\numberwithin{equation}{section}
\newcommand{\R}{\mathbb R}
\newcommand{\Z}{\mathbb Z}
\newcommand{\E}{\mathbb E}
\renewcommand{\P}{\mathbb P}
\newcommand{\ind}{\mathbf1}
\newcommand{\Rin}{R_{\mathrm{in}}}
\newcommand{\Rout}{R_{\mathrm{out}}}
\newcommand{\Rmom}{R_{\mathrm{mom}}}
\DeclareMathOperator{\dist}{dist}
\DeclareMathOperator{\supp}{supp}
\DeclareMathOperator{\Var}{Var}
\title[Inward-excited walks]{A general shape and fluctuation theorem for inward-excited walks}
\author{Ahmed Bou-Rabee}
\author{Tuan-Minh Nguyen}
\author{Yuval Peres}
\author{Arvind Singh}
\date{Four-author working draft, 13 October 2026}
\subjclass[2020]{Primary 60K35; Secondary 60G42, 60G50, 31C20}
\keywords{Excited random walk, convex limit shape, local times, martingale,
calibration, fluctuations}
\begin{document}
\begin{abstract}
We study nearest-neighbor walks whose first departure from each site has
an inward mean prescribed by a convex gauge, and whose later departures
are simple random walk. The first-departure laws may depend on the past,
and the drift strength may have spatial defects of codimension one.
One potential argument gives an almost-sure filled limit shape, uniform
conical local times, and recurrence. In the Euclidean case it retains
the sharp planar fluctuation exponent, polylogarithmic radius bounds
in higher dimensions, and the local-time limit laws. A quadratic
martingale gives Gaussian and iterated-logarithm laws for a weighted
radial moment for every convex body. An energy identity additionally
controls the number of sites with small local time. The proof combines
exit-time stopping, contact-point positivity, and a finite-scale
calibration, and includes the three previously studied models.
\end{abstract}
\maketitle

\section{The theorem and its specializations}

The range of an inward-excited walk is determined at first order by its
mean excitation. The probabilities realizing that mean, and spatial
changes confined to sufficiently thin sets, affect only lower-order
terms. We make this principle quantitative. The argument unifies
\citet{Singh2026}, \citet{BouRabeePeres2026CERW}, and
\citet{Nguyen2026}; all three models are special cases of the same
assumptions below.

Fix $d\ge2$. Let $K\subset\R^d$ be compact and convex, with $0$ in its
interior, and let $\rho(x)=\inf\{t>0:x\in tK\}$ be its gauge.
Write $V=|K|$ and $\omega_d=|B(0,1)|$. Symmetry and smoothness of $K$
are not required. Fix $\kappa>0$, subgradients
$\xi(x)\in\partial\rho(x)$ for $x\in\Z^d\setminus\{0\}$, and weights
$\lambda(x)\ge\lambda_->0$. Assume that, for a
constant $A<\infty$,
\begin{equation}\label{eq:model-defect}
 \sum_{x\in\Z^d\cap B(y,s)}|\lambda(x)-1|
 \le A s^{d-1}\qquad(y\in\R^d,\ s\ge1).
\end{equation}
Set $\xi(0)=0$ and $\lambda(0)=1$. Taking $s=1$ above gives
$\lambda(x)\le1+A$, so no separate upper bound is needed.

The walk starts at $X_0=0$. On its first departure from $x$, its
conditional step distribution, given the entire past, can be any
distribution on $\{\pm e_1,\ldots,\pm e_d\}$ with mean
\begin{equation}\label{eq:model-mean}
 -\kappa\lambda(x)\xi(x).
\end{equation}
On every later departure it chooses its neighbor uniformly. The
first departure from zero may have any mean-zero nearest-neighbor law.
Only existence of the distributions in \eqref{eq:model-mean} is
assumed; in particular no lower bound on individual probabilities is
required. The fields $\lambda$ and $\xi$ are deterministic. Allowing
the realizing distributions to depend on the past is useful because
every estimate below uses their conditional means and bounded steps.

Put
\begin{gather}
 \ell_n(x)=\sum_{j=0}^{n-1}\ind_{\{X_j=x\}},\qquad
 A_n=\{x:\ell_n(x)>0\},\qquad
 D_n=\bigcup_{x\in A_n}\bigl(x+[-1/2,1/2)^d\bigr),\label{eq:range-def}\\
 r_n=\left(\frac{(d+1)n}{2d\kappa V}\right)^{1/(d+1)},\quad
 b_n=\inf_{y\notin D_n}\rho(y),\quad
 B_n=\max_{j\le n}\rho(X_j).\label{eq:radius-def}
\end{gather}
Thus $\int 2d\kappa(r_n-\rho)_+=n$. Tildes denote constant extension
of lattice functions to the unit cells. All probability estimates
below concern integers $n\ge2$. Constants may depend on
$d,K,\kappa,\lambda_-,A$ and a specified probability exponent $p$,
but are uniform over the allowed first-departure laws.

\begin{theorem}[Shape, occupation, and moment universality]\label{thm:main}
Under \eqref{eq:model-defect}--\eqref{eq:model-mean}, the following hold.
\begin{enumerate}[label=\textup{(\roman*)},leftmargin=*]
\item Almost surely, for every $\varepsilon>0$ and all sufficiently
large $n$,
\[
 \{x\in\Z^d:\rho(x)<(1-\varepsilon)r_n\}
 \subset A_n\subset
 \{x\in\Z^d:\rho(x)<(1+\varepsilon)r_n\},
\]
and $r_n^{-1}\max_x|\ell_n(x)-2d\kappa(r_n-\rho(x))_+|\to0$.
In particular every site is visited infinitely often.

\item Define
\[
 h_n=\begin{cases}\sqrt{r_n\log n},&d=2,\\\log n,&d\ge3,
 \end{cases}\qquad q_n=h_n/r_n.
\]
For every $p>0$, with probability at least $1-C_pn^{-p}$,
\begin{align}
 |b_n-r_n|&\le C_pr_n\sqrt{q_n},\label{eq:main-inner}\\
 B_n-r_n&\le C_pr_nq_n^{1/(d+1)}\log n,\label{eq:main-outer}\\
 \max_x|\ell_n(x)-2d\kappa(r_n-\rho(x))_+|
 &\le C_pr_nq_n^{1/(d+1)}.\label{eq:main-profile}
\end{align}

\item Set
\[
 \delta_n=\begin{cases}\sqrt{r_n}\log n,&d=2,\\
 \sqrt{r_n\log n},&d\ge3.
 \end{cases}
\]
On an event of probability at least $1-C_pn^{-p}$, simultaneously
for every real $t\ge0$,
\begin{equation}\label{eq:main-occupation}
 \#\{x\in A_n:\ell_n(x)\le t\}
 \le C_p(t+\delta_n)r_n^{d-1}.
\end{equation}

\item Let
\[
 S_n=\sum_{x\in A_n}\lambda(x)\rho(x),\qquad
 J_K=\int_K|z|^2\,dz,\qquad
 \sigma_K^2=\frac{2J_K}{\kappa(d+3)}.
\]
Then
\begin{equation}\label{eq:main-moment-clt}
 \frac{S_n-n/(2\kappa)}{r_n^{(d+3)/2}}
 \Longrightarrow N(0,\sigma_K^2),
\end{equation}
and, for either sign, almost surely,
\begin{equation}\label{eq:main-moment-lil}
 \limsup_{n\to\infty}
 \frac{\pm(S_n-n/(2\kappa))}
 {r_n^{(d+3)/2}\sqrt{2\log\log n}}=\sigma_K.
\end{equation}
For $\Rmom(n)=((d+1)S_n/(dV))^{1/(d+1)}$, the same assertions
hold with $S_n-n/(2\kappa)$ replaced by $\Rmom(n)-r_n$,
the exponent $(d+3)/2$ replaced by $(3-d)/2$, and
$\sigma_K$ replaced by $\sigma_K/(dV)$.
\end{enumerate}
\end{theorem}

The weights in the moment statistic are essential: condition
\eqref{eq:model-defect} allows an $O(r_n^d)$ change in
$\sum_{A_n}\rho(x)$, which can exceed its Gaussian normalization.
They disappear when $\lambda\equiv1$.

For the Euclidean case put $\Rin=b_n$, $\Rout=B_n$, and
$W_n=\Rout-\Rin$. These definitions use cells for the inner radius
and include $X_n$ in the outer radius.

\begin{theorem}[Euclidean refinement]\label{thm:radial}
Under the assumptions of Theorem~\ref{thm:main}, with $K=\{|x|\le1\}$,
all the following conclusions hold for the same class of cookie laws.
\begin{enumerate}[label=\textup{(\roman*)},leftmargin=*]
\item For every $p>0$, with probability at least $1-C_pn^{-p}$,
\begin{align}
 |\Rin-r_n|&\le C_ph_n,\qquad
 |D_n\mathbin\triangle B(0,r_n)|\le C_pr_n^{d-1}h_n,
 \label{eq:radial-inner}\\
 \Rout-r_n&\le C_p\begin{cases}
 \sqrt{r_n}(\log n)^{5/2},&d=2,\\
 (\log n)^{d+1},&d\ge3,
 \end{cases}\label{eq:radial-outer}\\
 \max_x|\ell_n(x)-2d\kappa(r_n-|x|)_+|
 &\le C_p\begin{cases}
 \sqrt{r_n}(\log n)^{3/2},&d=2,\\
 \sqrt{r_n\log n},&d\ge3.
 \end{cases}\label{eq:radial-profile}
\end{align}
For each fixed $\theta<1$, restricting the maximum to
$|x|\le\theta r_n$ improves the planar bound to
$C_{p,\theta}\sqrt{r_n}\log n$.

\item Let $G$ be the simple-random-walk Green function when $d\ge3$,
and put
\[
 a_n=\begin{cases}\sqrt{r_n\log r_n},&d=2,\\
 \sqrt{r_n},&d\ge3,
 \end{cases}\quad
 \Sigma(y,z)=\begin{cases}16\kappa/\pi,&d=2,\\
 2d\kappa(2G(y-z)-\ind_{\{y=z\}}),&d\ge3.
 \end{cases}
\]
For any finite set of fixed sites, the vector
$a_n^{-1}(\ell_n(y)-2d\kappa(r_n-|y|))$ converges in distribution
to a centered Gaussian vector with covariance $\Sigma$.
At each fixed site and for either sign, almost surely,
\begin{equation}\label{eq:radial-site-lil}
 \limsup_n\frac{\pm(\ell_n(y)-2d\kappa(r_n-|y|))}
 {a_n\sqrt{2\log\log n}}=\sqrt{\Sigma(y,y)}.
\end{equation}

\item In dimension two, for every $p>0$ there is $c_p>0$ such that
each of
\begin{equation}\label{eq:radial-poly-lower}
 \P\{r_n-\Rin\ge c_p\sqrt{r_n\log n}\},\quad
 \P\{\Rout-r_n\ge c_p\sqrt{r_n\log n}\},\quad
 \P\{W_n\ge c_p\sqrt{r_n\log n}\}
\end{equation}
is at least $n^{-p}$ for all large $n$. Almost surely,
\begin{align}
 \limsup_n\frac{r_n-\Rin}{\sqrt{r_n\log\log n}},\quad
 \limsup_n\frac{\Rout-r_n}{\sqrt{r_n\log\log n}}
 &\ \ge\frac1{\sqrt{10\pi\kappa}},\label{eq:radial-lil-lower}\\
 \limsup_n\frac{W_n}{\sqrt{r_n\log\log n}}
 &\ge\sqrt{\frac\pi{3\kappa}}.\label{eq:radial-width-lil}
\end{align}
For every $a\ge0$,
\begin{equation}\label{eq:radial-smallball}
 \limsup_n\P\{W_n\le a\sqrt{r_n}\}
 \le1-\exp(-3\kappa a^2/\pi).
\end{equation}

\item In every dimension, for some $c>0$ and all large $n$,
\begin{equation}\label{eq:radial-local-lower}
 \P\left\{\max_{|x|\le\sqrt{r_n}}
 |\ell_n(x)-2d\kappa(r_n-|x|)|\ge
 c\begin{cases}\sqrt{r_n}\log n,&d=2,\\
 \sqrt{r_n\log n},&d\ge3
 \end{cases}\right\}\ge1-n^{-c}.
\end{equation}
\end{enumerate}
The bounds in part \textup{(i)} hold eventually almost surely.
\end{theorem}

We also retain the logarithmic radial lower bounds under the additional
assumption of a uniform positive lower bound on all first-departure
probabilities; their full tail forms are
\eqref{eq:lower-log-joint}--\eqref{eq:lower-log-outer} below.
The proof yields two additional boundary estimates. Write
$R=\Rout(n)+2\sqrt d$ and $E_n=D_n\setminus B(0,\Rin(n))$.
With the same probability as in Theorem~\ref{thm:radial}, uniformly
over unit vectors $u$ and $t\ge1$,
\begin{equation}\label{eq:new-caps}
 |E_n\cap B(Ru,t)|\le C_ph_n t^{d-1}.
\end{equation}
For lattice sites with $\Rin\le|y|\le R$, put $s=R-|y|+1$.
Then
\begin{equation}\label{eq:new-boundary}
 \ell_n(y)\le C_p\begin{cases}
 h_n+\sqrt{h_ns},&d=2,\\
 \log n+((\log n)s^{d-1})^{1/d},&d\ge3.
 \end{cases}
\end{equation}

\paragraph{Recovery of the three models.}
For the radial model of \citet{Singh2026}, take $K=\{|x|\le1\}$,
$\lambda\equiv1$, and $\kappa=\beta$. His prescribed-mean assumption
is exactly \eqref{eq:model-mean}. For the norm model of
\citet{BouRabeePeres2026CERW}, take $\lambda\equiv1$ and the particular
law $q_x(\pm e_i)=1/(2d)\mp\kappa\xi_i(x)/2$. Their strict parameter
restriction guarantees existence and, for the logarithmic lower
bounds, uniform ellipticity. Theorem~\ref{thm:radial} retains the
Euclidean constants and exponents, and
$J_{B(0,1)}=d\omega_d/(d+2)$ recovers their moment variances.

For \citet{Nguyen2026}, let $K=\{\|x\|_1\le1\}$,
$\xi(x)=\operatorname{sgn}(x)$, $\kappa=\theta/d$, and
$\lambda(x)=d/m(x)$ at $x\ne0$, where $m(x)$ is the number of its
nonzero coordinates. Use his first-departure kernel. The weights
differ from one only on the union of the coordinate hyperplanes,
which contains at most $Cs^{d-1}$ lattice points in any ball of radius
$s\ge1$. Thus \eqref{eq:model-defect} holds, including at $\theta=1$.
Theorem~\ref{thm:main} strengthens Hausdorff convergence in probability
to almost-sure filled inclusions and uniform local-time convergence.
Since $|K|=2^d/d!$, its radius is
$r_n=\tfrac12((d+1)!n/\theta)^{1/(d+1)}$, with the original
normalization. This verifies model containment at the transition-rule
level, including the exceptional coordinate hyperplanes.

\paragraph{The proof.}
Green martingales express local times as the potential of the occupied
set with a small error. Singh's exit-time stopping gives the initial occupation bound. At the first hole, convexity makes every excess contribution
nonnegative. This contact estimate and a single exponential crossing
bound give the shape. Bou-Rabee and Peres's contact and exterior-probe method gives the
fluctuation bounds. A finite-scale form of Nguyen's energy calibration supplies
\eqref{eq:main-occupation} directly. Once the cone is known, later departures dominate
all relevant quadratic variations, explaining the universality of the
Gaussian constants. We give the moderate-deviation argument separately
because a central limit theorem alone does not imply the lower tails.

\section{A common potential argument}

We use the assumptions and notation of Section~1. Put
$c=\min_{|u|=1}\rho(u)>0$, $\Lambda=\max_{|u|=1}\rho(u)$,
and $C_x=x+[-1/2,1/2)^d$. Write $q_x$ for the conditional law
used at the first departure from $x$; this law can be random,
but its prescribed mean is fixed. Every estimate below is pathwise
uniform over its possible values. Set $\Delta=P-I$ for the
simple-random-walk generator and
$\bar\nabla_i f(x)=(f(x+e_i)-f(x-e_i))/2$.
For a bounded measurable $D$, define
\begin{equation}\label{eq:core-potential}
 U_D(y)=\frac{2\kappa}{\omega_d}
 \int_D\nabla\rho(v)\cdot\frac{v-y}{|v-y|^d}\,dv.
\end{equation}

\begin{lemma}[Cone, layers, and contact]\label{lem:core-geometry}
For $b\geq0$,
$U_{bK}(y)=2d\kappa(b-\rho(y))_+$.
If $E\subset\{b\leq\rho\leq b+a\}$ then
$\|U_E\|_\infty\leq2d\kappa a$; for arbitrary $E$,
$\|U_E\|_\infty\leq C|E|^{1/d}$.
Suppose $bK\subset D\subset B(0,R)$ up to null sets, $y_0\in\partial(bK)$ lies in
$\overline{D^c}$, and $E=D\setminus bK$. Put
$H=U_D(y_0)$ and $I=\int_E(\rho-b)$.
Then
\begin{equation}\label{eq:core-contact}
 H\geq0,\qquad I\leq CR^d H,\qquad
 \|U_E\|_\infty\leq C I^{1/(d+1)}.
\end{equation}
\end{lemma}
\begin{proof}
For a compactly supported Lipschitz $f$, polar coordinates about $y$ give
$\int\nabla f(v)\cdot(v-y)|v-y|^{-d}\,dv=-d\omega_df(y)$.
Apply this to $f=(b-\rho)_+$. On each ray $y+tu$, convexity makes
$t\mapsto\rho(y+tu)$ decrease and then increase. On either part, its total
variation while $b\leq\rho\leq b+a$ is at most $a$; the two signed parts
therefore give an integral between $-a$ and $a$. Integrating over directions
proves the layer bound. Rearrangement of $|v-y|^{1-d}$ proves the volume
bound. Finally
$\nabla\rho(v)\cdot(v-y_0)\geq\rho(v)-\rho(y_0)=\rho(v)-b$
on $E$, proving the first two assertions in \eqref{eq:core-contact}.
Split $E$ at $\rho=b+a$: the layer contributes at most $Ca$ and the rest,
of volume at most $I/a$, contributes at most $C(I/a)^{1/d}$.
Take $a=I^{1/(d+1)}$; the case $I=0$ follows directly.
\end{proof}

\begin{lemma}[Discrete potential representation]\label{lem:core-green}
Let $g$ be the potential kernel in dimension two and minus the Green
function in dimensions at least three, normalized by
$\Delta g=\mathbf1_{\{0\}}$, where $\Delta=P-I$ for simple random walk.
Write $g_y(x)=g(x-y)$ and $M^y=M^{g_y}$ for its Dynkin martingale;
for finite combinations of these kernels, extend $M^f$ linearly.
These martingales have bounded increments. Whenever
$\max_{j\leq t}|X_j|\leq R$, $R\geq2$, and $y\in\Z^d$ satisfies $|y|\leq R^{d+2}$,
\begin{equation}\label{eq:core-green}
 \ell_t(y)=U_{D_t}(y)-M_t^y+O(\log R),\qquad
 \langle M^y\rangle_t\leq C W_t(y),\quad
 W_t(y)=\sum_x\ell_t(x)(1+|x-y|)^{2-2d}.
\end{equation}
Moreover, for $|z-y|\leq\sqrt d$,
$|U_{D_t}(z)-U_{D_t}(y)|\leq C\log R$.
\end{lemma}
\begin{proof}
Standard Green estimates \citep{LawlerLimic2010} give
\begin{align*}
 \bar\nabla g(z)&=\frac2{\omega_d}\frac z{|z|^d}+O((1+|z|)^{-d})
 \quad(z\ne0),\\
 |g(z+e)-g(z)|&\leq C(1+|z|)^{1-d},\\
 |g(z+e_i)+g(z-e_i)-2g(z)|&\leq C(1+|z|)^{-d}.
\end{align*}
Write $s_i(x)=q_x(e_i)+q_x(-e_i)$. Exactly,
\[
 (Q_x-P)g_y(x)=-\kappa\lambda(x)\xi(x)\cdot\bar\nabla g(x-y)
 +\frac12\sum_i(s_i(x)-1/d)
 [g(x+e_i-y)+g(x-e_i-y)-2g(x-y)].
\]
Dynkin's formula, grouping first departures by site, gives the first
identity in \eqref{eq:core-green} with the discrete drift potential.
The displayed second differences sum to $O(\log R)$. The endpoint
$g(X_t-y)-g(-y)$ has this bound too. By dyadic annuli,
\begin{equation}\label{eq:core-defect-potential}
 \sup_y\sum_{x\in B(0,CR)}
 |\lambda(x)-1|(1+|x-y|)^{1-d}\leq C\log R.
\end{equation}
Thus replacing $\lambda$ by one costs $O(\log R)$, even for an arbitrary
subset of visited defect sites. Replacing the Green gradient by its
leading term, and then by cell averages, has the same cost.

For completeness the subgradient replacement also costs only $O(\log R)$,
including at corners of $K$. The positive measure $\mu=\Delta_{\rm cont}\rho$
satisfies $\mu(B(y,s))\leq Cs^{d-1}$, by testing against a cutoff and using
$|\nabla\rho|\leq\Lambda$. The convex-cell estimate is
\begin{equation}\label{eq:core-cell}
 \int_{C_x}|\nabla\rho(v)-\xi(x)|\,dv
 \leq C\mu(B(x,6\sqrt d)).
\end{equation}
Indeed, $h(v)=\rho(v)-\rho(x)-\xi(x)\cdot(v-x)$ is nonnegative, convex,
and zero at $x$. Convexity along coordinate lines gives
$|\partial_i h(v)|\leq h(v+e_i)+h(v-e_i)$, so the left side is at most
$C\int_{B(x,2\sqrt d)}h$. Mollify $h$ and put $A=2\sqrt d$. Convexity gives
\[
 h_\varepsilon(x+tu)\le h_\varepsilon(x)
       +t\,\partial_r h_\varepsilon(x+2Au)
       \qquad(0\le t\le A,\ |u|=1).
\]
Integration in $t,u$ and the divergence theorem bound the integral by
$C h_\varepsilon(x)+C\mu(B(x,6\sqrt d))$.
Letting $\varepsilon\downarrow0$ proves \eqref{eq:core-cell}.
Bound the cells close to $y$ directly; bounded overlap and dyadic annuli
bound the remaining error by
$C\int_{1\leq|z-y|\leq CR^{d+2}}(1+|z-y|)^{1-d}\,d\mu(z)=O(\log R)$.
The increment estimate gives the bracket bound. Splitting the integral
at distance $2\sqrt d$ and differentiating the far kernel gives the
cell modulus.
\end{proof}

\begin{lemma}[Exit bootstrap]\label{lem:core-bootstrap}
Put $\tau_R=\inf\{t:\rho(X_t)\geq R\}$. For every $p>0$, outside an event
of probability $CR^{-p}$, simultaneously for $t\leq\tau_R$,
\begin{equation}\label{eq:core-bootstrap}
 \max_x\ell_t(x)\leq CR,\qquad
 \tau_R\leq CR^{d+1},\qquad
 \sup_{|y|\leq R^{d+2}}|M_t^y|\leq C\sqrt R\log R.
\end{equation}
Consequently
$\|\widetilde\ell_t-U_{D_t}\|_\infty\leq C\sqrt R\log R$ on
$B(0,CR)$.
\end{lemma}
\begin{proof}
To see that the exit time is finite, stop the compensated coordinate
martingale $Z_t^i=X_t^i+\kappa\sum_{x\in A_t}\lambda(x)\xi_i(x)$
at $\tau_R$. It is bounded by a deterministic constant $C_R$.
There are at most $m_R=O(R^d)$ first departures before exit, so its
bracket is at least $((t\wedge\tau_R)-m_R)/d$. Consequently
$\E(t\wedge\tau_R)\le m_R+dC_R^2$, and $\tau_R<\infty$ almost surely. Before exit, all departed sites lie in $B(0,CR)$.
The exact discrete Dynkin identity and the first Green increment bound
give $\ell_t(y)\leq |M_t^y|+CR$ without any preliminary clock estimate.
Stop also when one of the displayed martingales first reaches $R$ in
absolute value. Up to this stopping time, $\max_x\ell(x)\leq C'R$,
and hence the clock is at most $C'R^{d+1}$ and every bracket at most
$C'R\log R$. Freedman's maximal inequality \citep[Theorem~1.6]{Freedman1975} and a union bound over the
polynomially many poles give
$\sup|M^y|\leq C_p\sqrt R\log R$ with probability $1-CR^{-p}$.
This is smaller than $R$ for large $R$, so the additional stop is never
reached before $\tau_R$. The representation and cell modulus finish.
\end{proof}

\begin{lemma}[Exponential crossing and projection]\label{lem:core-crossing}
For every $p>0$ and integer $N\geq2$, outside an event of probability
$CN^{-p}$, simultaneously for $t\leq N$ and integers $0\leq k\leq C_0N$,
\begin{equation}\label{eq:core-crossing}
 \max_{j\leq t}\operatorname{dist}(X_j,kK)
 \leq C\left(1+\max_{\rho(x)\geq k}\ell_t(x)\right)\log N.
\end{equation}
Here $C_0$ is any fixed constant.
\end{lemma}
\begin{proof}
We first record the short crossing estimate. Fix a unit vector $u$, $a\geq0$,
and suppose all departed sites in the cap $u\cdot x>a$ have
$u\cdot\lambda(x)\xi(x)\geq\alpha>0$. With
$f(x)=e^{\theta u\cdot x}$, conditional increments satisfy
\[
 \mathbb E[f(X_{j+1})-f(X_j)\mid\mathcal F_j]
 \leq f(X_j)(C\theta^2-\kappa\alpha\theta I_j)
 \quad (u\cdot X_j>a),
\]
where $I_j$ indicates a first departure. Stop at time $N$, on entering
a cap site violating the drift condition, when a cap local time exceeds
$L$, or when $u\cdot X$ reaches $a+h$. Choose $\theta=c/(L+1)$.
Summing by site makes the total contribution from the cap nonpositive:
each departed site supplies its cookie once and at most $L+1$ visits.
The other at most $N$ increments contribute at most
$CN e^{\theta(a+1)}$. Optional stopping of the Dynkin martingale gives
\begin{equation}\label{eq:core-exp-crossing}
 \mathbb P\left(\begin{array}{c}
  \text{the level }a+h\text{ is reached by }N,\text{ all previously visited}\\
  \text{cap sites satisfy the drift condition, and cap local times}\leq L
 \end{array}\right)
 \leq CN\exp\{-ch/(L+1)\}.
\end{equation}
The stopping argument does not condition on the event in this display.

Let $w=\max_{j\leq t}\operatorname{dist}(X_j,kK)$ and choose $x_*=X_j$
attaining it. When $w>0$, let $\pi$ denote Euclidean projection onto $kK$,
$u=(x_*-\pi(x_*))/w$, and $\beta=u\cdot\pi(x_*)$.
Then $u\cdot X_j\leq T:=\beta+w$ for all $j\leq t$ and $\beta\geq0$.
Put $\eta=c^2/(8\Lambda^2)$. If $x$ is a site visited by time $t$
and $u\cdot x>T-\eta w$, write $v=(x-\pi(x))/|x-\pi(x)|$.
The projection inequalities give $u\cdot v>1-\eta$; moreover
\[
 \xi(x)\cdot v\geq
 \frac{\rho(x)-k}{\operatorname{dist}(x,kK)}\geq c,
\]
where the last inequality follows by comparing $\pi(x)$ with
$kx/\rho(x)$. Thus $u\cdot\lambda(x)\xi(x)\geq\lambda_-c/2$.
All these cap sites satisfy $\rho(x)>k$.
Use \eqref{eq:core-exp-crossing} with $a$ the least integer at least
$T-\eta w$, and with $L$ the integer above the largest relevant local time.
The directions $u$ belong to the deterministic list formed by projecting
the lattice sites $|x|\leq N$ onto $kK$, for integer $k\leq C_0N$.
This list, the integer levels $a\leq CN$, and $1\leq L\leq N+1$ have
polynomial total size. For each triple choose the deterministic length
$C_p(L+1)\log N$ in \eqref{eq:core-exp-crossing}. A union bound excludes
all these crossings. The drift condition is needed only at visited
sites, since the stopped estimate ends on entering any violating site. Since $T-a\geq\eta w-1$, this proves
\eqref{eq:core-crossing}; bounded $w$ is absorbed into the constant.
\end{proof}

\begin{proposition}[Coarse scale, with polynomial tails]\label{prop:core-coarse}
For every $p>0$ there are $c_p,C_p>0$ such that, with probability at least
$1-C_pR^{-p}$,
\begin{equation}\label{eq:core-clock}
 c_pR^{d+1}\leq\tau_R\leq C_pR^{d+1}.
\end{equation}
Consequently, with $r_n=((d+1)n/(2d\kappa|K|))^{1/(d+1)}$, with probability
at least $1-C_p n^{-p}$,
\begin{equation}\label{eq:core-coarse}
 \max_{j\leq n}|X_j|\leq C_pr_n,\quad
 c_pr_n\leq\max_x\ell_n(x)\leq C_pr_n,\quad
 c_pr_n^d\leq|A_n|\leq C_pr_n^d.
\end{equation}
\end{proposition}
\begin{proof}
Work on the events of the preceding two lemmas with a deterministic
horizon $N=\lceil CR^{d+1}\rceil$. At time $\tau_R$, let
$b=\inf_{y\notin D_{\tau_R}}\rho(y)$, and choose a contact point
$y_0\in\partial(bK)\cap\overline{D_{\tau_R}^c}$ in the closure of an
unoccupied cell. Continuity and the cell estimate imply
$0\leq H=U_{D_{\tau_R}}(y_0)\leq\delta_R=C\sqrt R\log R$.
Lemma~\ref{lem:core-geometry} therefore gives
\[
 \|\widetilde\ell_{\tau_R}-2d\kappa(b-\rho)_+\|_\infty
 \leq\delta_R+C(R^d\delta_R)^{1/(d+1)}
 =:\varepsilon_R=o(R/\log R).
\]
Use \eqref{eq:core-crossing} with $k=\lfloor b\rfloor$; outside $kK$ the
local times are at most $C+\varepsilon_R$. Since
$\rho(X_{\tau_R})\geq R$ and $\rho$ is $\Lambda$-Lipschitz,
\[
 R-b\leq C+C(\varepsilon_R+1)\log R=o(R).
\]
The cone is supported on $\{\rho<b\}\subset D_{\tau_R}$.
Integrating its approximation over $D_{\tau_R}\subset B(0,CR)$ gives
\[
 \tau_R=\frac{2d\kappa|K|}{d+1}b^{d+1}
 +O(R^d\varepsilon_R)\asymp R^{d+1},
\]
proving \eqref{eq:core-clock}. To obtain \eqref{eq:core-coarse}, choose
$R=\lceil C n^{1/(d+1)}\rceil$ so that the lower clock bound exceeds $n$,
and apply the exit event with exponent $(d+1)(p+1)$. Then $n<\tau_R$,
so the spatial and local-time upper bounds follow. The remaining bounds
follow from $n=\sum_x\ell_n(x)\leq|A_n|\max_x\ell_n(x)$ and
$|A_n|\leq C(1+\max_{j\leq n}|X_j|)^d$.
\end{proof}

\section{Quantitative stability of the cone}

Fix $p>0$, write $r=r_n$, $a=\log n$, and work on one event of
probability at least $1-Cn^{-p}$. Proposition~\ref{prop:core-coarse},
Lemma~\ref{lem:core-green}, Freedman's inequality and a union bound give,
simultaneously for all lattice sites $|y|\le3n$,
\begin{equation}\label{eq:rate-event}
 D_n\subset B(0,Cr),\quad \max\ell_n\le Cr,\quad
 |M_n^y|\le C(\sqrt{aW_n(y)}+a),\quad
 \ell_n(y)=U_{D_n}(y)-M_n^y+O(a).
\end{equation}
Here and below the constants may be increased by increasing the exponent
in the constituent probability estimates. Cell interpolation gives
\begin{equation}\label{eq:rate-crude}
 \|\ell_n-U_{D_n}\|_\infty\le\delta,
 \qquad \delta=
 \begin{cases}C\sqrt r\,a,&d=2,\\ C\sqrt{ra},&d\ge3.
 \end{cases}
\end{equation}
For points outside $B(0,2n)$ this follows directly from the decay of the
potential, since $\ell_n=0$ there. Define
\[
 b=\inf_{y\notin D_n}\rho(y),\quad E=D_n\setminus bK,
 \quad H=U_{D_n}(y_0),\quad I=\int_E(\rho(v)-b)\,dv,
\]
where $y_0\in\partial(bK)$ is in the closure of an empty cell $C_z$.
Thus $\ell_n(z)=0$, $|z-y_0|\le\sqrt d/2$, and convexity gives
$H=U_E(y_0)\ge0$. Set
\begin{equation}\label{eq:rate-hq}
 h=\begin{cases}\sqrt{ra},&d=2,\\a,&d\ge3,\end{cases}
 \qquad q=h/r.
\end{equation}

\begin{lemma}[Noise at the first hole]\label{lem:rate-contact}
On the preceding event, $H\le Ch$.
\end{lemma}
\begin{proof}
The equation at the empty site and the cell modulus give
\begin{equation}\label{eq:rate-contact-noise}
 H\le C(\sqrt{aW_n(z)}+a).
\end{equation}
For every nonnegative integrable radial function $\chi$, Newton's shell
identity gives, with $m(s)=\int_{|v|<s}\chi(v)\,dv$,
\begin{equation}\label{eq:rate-newton}
 (\chi*U_D)(y)=\frac{2\kappa}{\omega_d}\int_D
 \nabla\rho(v)\cdot\frac{v-y}{|v-y|^d}m(|v-y|)\,dv.
\end{equation}
Indeed the spherical average of $(v-y-w)/|v-y-w|^d$ on $|w|=s$
is $(v-y)|v-y|^{-d}\mathbf1_{\{|v-y|>s\}}$; integrate in $s$.
Since the contact integrand is nonnegative,
$(\chi*U_E)(y_0)\le\|\chi\|_1H$.

In the plane take $\chi(v)=(1+|v|)^{-2}\mathbf1_{\{|v|\le Cr\}}$,
with the cutoff large enough to cover $D_n-y_0$.
Cell comparison and \eqref{eq:rate-crude} give
\[
 W_n(z)\le C(\chi*\ell_n)(y_0)
 \le C(\chi*U_{bK})(y_0)+C(H+\delta)\|\chi\|_1.
\]
The cone $U_{bK}=2d\kappa(b-\rho)_+$ is Lipschitz and vanishes at
$y_0$, so its convolution at $y_0$ is $O(r)$. Also
$\|\chi\|_1=O(a)$ and $H\le\delta$ by \eqref{eq:rate-crude}.
Consequently $W_n(z)\le C(r+\sqrt r\,a^2)\le Cr$, proving the claim.

For $d\ge3$ put $\chi(v)=(1+|v|)^{2-2d}$. This kernel is integrable
and $\chi*\chi\le A\chi$: split the convolution into
$|w|\ge|v|/2$ and $|v-w|\ge|v|/2$. Cell comparison, the martingale
bound and Young's inequality yield, for every fixed $\gamma>0$,
\[
 \ell_n\le U_{D_n}+C_\gamma a+\gamma\chi*\ell_n.
\]
This inequality holds globally; outside $B(0,2n)$ the local time is zero
and $|U_{D_n}|\le C$. Convolve it with $\chi$ and take
$\gamma=(2A)^{-1}$ to obtain the resolvent estimate
\begin{equation}\label{eq:rate-resolvent}
 \chi*\ell_n\le2\chi*U_{D_n}+Ca.
\end{equation}
The cone is Lipschitz, bounded by $Cb$, and $b\le Cr$, hence
\begin{equation}\label{eq:rate-cone-convolution}
 (\chi*U_{bK})(y)\le\|\chi\|_1 U_{bK}(y)+Ca;
\end{equation}
the error is bounded by $C\int\chi(v)\min\{|v|,r\}\,dv=O(a)$.
At $y_0$, \eqref{eq:rate-newton} and \eqref{eq:rate-resolvent} now give
$W_n(z)\le C(\chi*\ell_n)(y_0)\le C(H+a)$.
Substitute in \eqref{eq:rate-contact-noise} and absorb
$C\sqrt{aH}\le H/2+Ca$.
\end{proof}

\paragraph{Rates for every convex body.}
Lemma~\ref{lem:core-geometry} and the contact bound give
$I\le Cr^d h=Cr^{d+1}q$. Split $E$ into the layer
$b\le\rho<b+s$, whose volume is at most $Cr^{d-1}s$, and the remainder,
whose volume is at most $I/s$. With $s=r\sqrt q$ this proves
\begin{equation}\label{eq:rate-general-mass}
 |E|\le Cr^d\sqrt q,\qquad
 \|U_E\|_\infty\le Crq^{1/(d+1)}.
\end{equation}
The quadratic martingale is
\[
 Q_n=|X_n|^2-n+2\kappa\sum_{x\in A_n}\lambda(x)\rho(x).
\]
Stopping on leaving $B(0,Cr)$ and applying Azuma gives
$|Q_n|\le Cr\sqrt{na}$. The sparse-weight condition and cell quadrature
each cost $O(r^d)$ in replacing the weighted sum by $\int_{D_n}\rho$.
Since $\int_{bK}\rho=d|K|b^{d+1}/(d+1)$ and
$n=2d\kappa|K|r^{d+1}/(d+1)$, we obtain
\begin{equation}\label{eq:rate-clock-comparison}
 \left|1-(b/r)^{d+1}\right|
 \le C\left(r^{(1-d)/2}\sqrt a+r^{-1}+|E|r^{-d}\right).
\end{equation}
The first two terms are at most $C\sqrt q$; hence
\begin{equation}\label{eq:rate-general-conclusion}
 |b-r|\le Cr\sqrt q,\qquad
 \max_x|\ell_n(x)-2d\kappa(r-\rho(x))_+|
 \le Crq^{1/(d+1)}.
\end{equation}
Here we used the cone identity, \eqref{eq:rate-crude} and
\eqref{eq:rate-general-mass}; their errors are bounded by the right side.
Apply the projection crossing estimate \eqref{eq:core-crossing} with
$k=\lfloor r\rfloor$. Outside $kK$ the local time is at most
$C(1+rq^{1/(d+1)})$. Since $\rho$ is Lipschitz, it follows that
\begin{equation}\label{eq:rate-general-outer}
 \max_{j\le n}\rho(X_j)\le r+Crq^{1/(d+1)}a.
\end{equation}
These prove Theorem~\ref{thm:main}\textup{(ii)}. Taking $p>1$ and
applying Borel--Cantelli gives the filled inclusions and uniform cone
law in part \textup{(i)}; at each fixed site the cone grows linearly in
$r_n$, proving recurrence.

\section{Euclidean geometry and exterior probes}
Now let $\rho(v)=|v|$. We retain $r,a,b,E,h,q$ above, so $b=R_{\rm in}(n)$.
For $|v|\ge b=|y_0|$ the exact identity
\[
 2(|v|^2-v\cdot y_0)=|v-y_0|^2+|v|^2-b^2
\]
gives
\[
 u_v\cdot\frac{v-y_0}{|v-y_0|^d}
 \ge\frac{|v-y_0|^{2-d}}{2|v|}\ge c r^{1-d}
 \qquad(v\in E).
\]
Thus the contact estimate improves the volume and inner radius bounds to
\begin{equation}\label{eq:rate-euclidean-inner}
 |E|\le Cr^{d-1}h=Cr^dq,\quad |b-r|\le Crq,
 \quad |D_n\mathbin\triangle B(0,r)|\le Cr^dq.
\end{equation}
The radius assertion follows from \eqref{eq:rate-clock-comparison} since
its first two terms are at most $Cq$.

Write $U_E^+$ for the radial potential with its negative integrand
replaced by zero, and put
\[
 L=1+\max_{|x|\ge b}\ell_n(x),\qquad
 w=R_{\rm out}(n)-b+2\sqrt d,\qquad R=b+w.
\]
Then $w\ge1$ and $E\subset\{b\le|v|\le R\}$.
Crossing with $k=\lceil b\rceil$ gives $w\le CLa$.
For $d\ge3$, \eqref{eq:rate-newton} gives
\[
 \chi*U_E\le\|\chi\|_1U_E^+,\qquad
 \chi*U_E\le\|\chi\|_1\|U_E\|_\infty.
\]
Combining these with the resolvent and cone-convolution estimates yields
\begin{equation}\label{eq:rate-envelope}
 \ell_n(y)\le C\bigl((b-|y|)_++\|U_E\|_\infty+a\bigr),
 \quad
 \ell_n(y)\le C(U_E^+(y)+a)\quad(|y|\ge b).
\end{equation}
In the plane use \eqref{eq:rate-crude} instead. In either case the
volume bound $\|U_E\|_\infty\le C|E|^{1/d}$ gives
\begin{equation}\label{eq:rate-preliminary-width}
 L\le Crq^{1/d},\qquad w\le Crq^{1/d}a=o(r a^{-d}).
\end{equation}
In particular $b\sim r$ and $w=o(r)$. The layer estimate also gives
$\|U_E\|_\infty\le Cw$.

\paragraph{One exterior estimate in all dimensions.}
Set
\begin{equation}\label{eq:rate-theta}
 \Theta=\begin{cases}
 h,&d=2,\\
 a+r^{-1}\displaystyle\sup_y\int_E|v-y|^{2-d}\,dv,&d\ge3.
 \end{cases}
\end{equation}
In either case $\Theta\ge a$ for large $n$ and
$\sup_y\int_E|v-y|^{2-d}\,dv\le Cr\Theta$; in the plane this is
just $|E|\le Crh$. For $|y|\ge R$, the Euclidean identity therefore gives
$U_E^+(y)\le Cr^{-1}\int_E|v-y|^{2-d}\,dv\le C\Theta$.

The planar substitute for the resolvent is particularly short.
The layer estimate and \eqref{eq:rate-crude} give
$\ell_n(x)\le4\kappa(b-|x|)_++C(w+\sqrt r\,a)$.
For every lattice $y$ with $|y|\ge b$, inserting this in $W_n(y)$ gives
\[
 W_n(y)\le Cr+C(w+\sqrt r\,a)a\le Cr.
\]
Indeed $(b-|x|)_+\le|x-y|$ bounds the cone contribution by $Cr$,
and \eqref{eq:rate-preliminary-width} makes the other term $o(r)$.
Thus the martingale bound, and decay beyond the pole set, imply
\begin{equation}\label{eq:rate-planar-exterior}
 |\ell_n(y)-U_E(y)|\le Ch\qquad(d=2,\ |y|\ge b).
\end{equation}
Together with \eqref{eq:rate-envelope}, this gives in every dimension
$\ell_n(y)\le C(U_E^+(y)+\Theta)$ outside the inner ball.

We now use the same exterior probe in every dimension. Given a unit
vector $u$ and $\sqrt d\le t\le r/4$, put $y_t=(R+4t)u$ and let $z$
be its cell center. Since $R=R_{\rm out}+2\sqrt d$, this cell is empty
and $|z|\ge R$. If $d\ge3$, the resolvent, cone-convolution estimate
and $\chi*U_E\le\|\chi\|_1U_E^+$ give $W_n(z)\le C\Theta$;
if $d=2$, use \eqref{eq:rate-planar-exterior}.
The martingale estimate and cell modulus give in both cases
\begin{equation}\label{eq:rate-probe}
 |U_E(y_t)|\le C\Theta.
\end{equation}
For $v\in E\cap B(Ru,t)$ we have $u_v\cdot u\ge1-t/R$ and
\[
 u_v\cdot(v-y_t)\le R-(R+4t)(1-t/R)\le-t,
 \qquad |v-y_t|\le5t,
\]
because $t\le R/2$ for large $n$. Such points contribute at most
$-ct^{1-d}$ each, whereas all positive contributions total at most
$C\Theta$. Hence \eqref{eq:rate-probe} implies
\begin{equation}\label{eq:rate-cap}
 |E\cap B(Ru,t)|\le C\Theta t^{d-1}\qquad(|u|=1,\ t\ge1).
\end{equation}
For $t>r/4$ use $|E|\le Cr^{d-1}h\le Cr^{d-1}\Theta$;
for $1\le t<\sqrt d$ use the ball-volume bound. Thus the estimate
holds down to a fixed spatial scale, independently of the shell width.

\begin{lemma}[Cap volume controls the boundary profile]\label{lem:rate-near-far}
For every $y$ with $b\le|y|\le R$, put $s=R-|y|+1$. Then
\[
 U_E^+(y)\le C(\Theta s^{d-1})^{1/d}+C\Theta.
\]
The same bound holds for $\ell_n(y)$ at lattice sites.
\end{lemma}
\begin{proof}
For $t\ge s$, $B(y,t)\subset B(Ru_y,2t)$, so
\eqref{eq:rate-cap} gives $|E\cap B(y,t)|\le C\Theta t^{d-1}$.
The part $|v-y|\le s$ has volume at most $C\Theta s^{d-1}$ and
therefore potential at most $C(\Theta s^{d-1})^{1/d}$.
For the rest, $|v|^2-|y|^2\le2Rs$ and $|v|\ge b\asymp r$ give
\[
 \left[u_v\cdot\frac{v-y}{|v-y|^d}\right]_+
 \le\frac{Cs}{|v-y|^d}+\frac{C}{r|v-y|^{d-2}}.
\]
The cap bound makes the first integral at most
$Cs\int_s^\infty\Theta t^{d-1}t^{-d-1}\,dt\le C\Theta$;
the definition of $\Theta$ bounds the second by $C\Theta$.
The exterior local-time bound proves the last assertion.
\end{proof}

Since $s\le w+1\le2w$, the lemma and crossing give
\begin{equation}\label{eq:rate-L}
 L\le C\bigl((\Theta w^{d-1})^{1/d}+\Theta\bigr),
 \qquad w\le CLa\le\tfrac12w+C\Theta a^d.
\end{equation}
Thus $w\le C\Theta a^d$ in every dimension. In the plane $\Theta=h$,
so the lemma immediately yields
\begin{equation}\label{eq:rate-planar-final}
 w\le C\sqrt r\,a^{5/2},\qquad
 \sup_{b\le|y|\le R}U_E^+(y)\le C\sqrt r\,a^{3/2}.
\end{equation}

For $d\ge3$ it remains only to absorb the curvature term in $\Theta$.
If $t\ge w$ and $B(y,t)$ meets $E$ at $v$, its intersection with $E$
lies in $B(v,2t)$ and hence in a cap of radius at most $3t$.
Consequently
\[
 |E\cap B(y,t)|\le C\min\{\Theta t^{d-1},r^{d-1}a\}\quad(t\ge w),
 \qquad |E\cap B(y,t)|\le\omega_dt^d\quad(t<w).
\]
Layer-cake integration, splitting the minimum at
$t=r(a/\Theta)^{1/(d-1)}$, gives
\[
 \sup_y\int_E|v-y|^{2-d}\,dv
 \le Cw^2+Cr\Theta^{(d-2)/(d-1)}a^{1/(d-1)}.
\]
Substitution in \eqref{eq:rate-theta} and Young's inequality give
\begin{equation}\label{eq:rate-theta-improved}
 \Theta\le C(a+w^2/r).
\end{equation}
Together with $w\le C\Theta a^d$, this implies
$w\le Ca^{d+1}+Ca^dw^2/r$. The last term is at most $w/2$ for
large $n$ by \eqref{eq:rate-preliminary-width}. We conclude that
\begin{equation}\label{eq:rate-higher-final}
 w\le Ca^{d+1},\qquad \Theta\le Ca.
\end{equation}
Substitution in \eqref{eq:rate-cap} proves the all-scale cap estimates.
Lemma~\ref{lem:rate-near-far} now also gives
\begin{equation}\label{eq:rate-boundary-profile}
 \ell_n(y)\le
 \begin{cases}
 C\bigl(h+\sqrt{hs}\bigr),&d=2,\\
 C\bigl(a+(as^{d-1})^{1/d}\bigr),&d\ge3,
 \end{cases}
 \qquad b\le|y|\le R,\quad s=R-|y|+1.
\end{equation}
In particular, for $d\ge3$ the outermost radial layer of width
$C_0\log n$ has local times $O(\log n)$. If this layer reaches inside
$B(0,b)$, then $w=O(a)$. The continuum boundary bound from the lemma
and the maximum principle give $U_E\le Ca$ inside that ball.
There $\chi*U_E\le\|\chi\|_1U_E$, since the potential integrand is
nonnegative; the resolvent and the cone depth $O(a)$ give the same
local-time bound.

\paragraph{Uniform and interior local times.}
For $d\ge3$, the layer bound and \eqref{eq:rate-higher-final} give
$\|U_E\|_\infty\le Ca^{d+1}$. Together with the cone identity,
\eqref{eq:rate-crude} and \eqref{eq:rate-euclidean-inner}, this gives
\[
 \max_x|\ell_n(x)-2d\kappa(r-|x|)_+|\le C\sqrt{ra}.
\]
In the plane $U_E$ is nonnegative and harmonic on $B(0,b)$, and continuous
on its closure. The maximum principle and \eqref{eq:rate-planar-final}
therefore bound it throughout that ball by $C\sqrt r\,a^{3/2}$.
On the outer annulus use \eqref{eq:rate-planar-exterior}; beyond the range
both the local time and, except for a radial error $C|b-r|$, the target
cone vanish. Consequently the uniform planar error is
$C\sqrt r\,a^{3/2}$.

Finally, for any fixed $\theta<1$ and $|y|\le\theta r$, the excess set
is at distance at least $c_\theta r$ for large $n$. Thus
\begin{equation}\label{eq:rate-bulk-potential}
 |U_{D_n}(y)-2d\kappa(r-|y|)|
 \le C_\theta(r^{1-d}|E|+|b-r|)\le C_\theta rq.
\end{equation}
The corresponding local-time error is at most $C_\theta\sqrt r\,a$
in the plane and $C_\theta\sqrt{ra}$ in higher dimensions.
All assertions hold simultaneously with probability $1-Cn^{-p}$ and,
by Borel--Cantelli, eventually almost surely. The all-scale cap bounds
\eqref{eq:rate-cap} and boundary profile \eqref{eq:rate-boundary-profile}
are additional quantitative conclusions of the same proof.

\section{Calibration and small occupation layers}
\begin{lemma}[Finite-scale calibration]\label{lem:finite-calibration}
Let $a=2d\kappa$, $b=\nabla\rho$, and let $D\subset B(0,R)$ be measurable,
$R\geq1$. Put $u=U_D$ and suppose
\[
 u\geq-\delta\quad\hbox{on }B(0,2R),\qquad
 |u|\leq\delta\quad\hbox{on }B(0,2R)\setminus D.
\]
Then $w=(u-\delta)_+$ vanishes off $D$, belongs to $H^1(\mathbb R^d)$,
and, with $A=\{u>\delta\}$,
\begin{equation}\label{eq:finite-calibration}
 \int_{\mathbb R^d}|\nabla w+a\mathbf1_A b|^2
 +a^2\int_{D\setminus A}|b|^2
 \leq C\delta R^{d-1}.
\end{equation}
In particular $|D\setminus A|\leq C\delta R^{d-1}$.
\end{lemma}
\begin{proof}
The potential is continuous, tends to zero at infinity, and is harmonic
off $\overline D$. The maximum principle extends $|u|\leq\delta$ from the
annulus to the exterior of $B(0,R)$. The potential equation is
$\operatorname{div}(\nabla u+a\mathbf1_D b)=0$.
Since $\mathbf1_D b\in L^2$, the usual Fourier multiplier estimate gives
$\nabla u\in L^2$; hence the compactly supported truncation $w$ is in $H^1$.
Testing the equation with $w$, using $A\subset D$, gives
\begin{equation}\label{eq:calibration-test-w}
 \int|\nabla w|^2=-a\int b\cdot\nabla w.
\end{equation}
Choose a smooth $\chi$ equal to one on $B(0,R)$, zero off $B(0,2R)$,
with $|\nabla\chi|\leq C/R$ and $|\Delta\chi|\leq C/R^2$.
Put $v=u-w$, so $|v|\leq\delta$. The equation tested with $\chi\rho$,
together with \eqref{eq:calibration-test-w}, gives the exact identity
\[
 \int|\nabla w+a\mathbf1_D b|^2
 =a\int b\cdot\nabla w+a^2\int_D|b|^2
 =-a\int\nabla v\cdot\nabla(\chi\rho)
 =a\int v\,d\Delta(\chi\rho).
\]
The convexity measure $\mu=\Delta\rho$ has
$\mu(B(0,2R))\leq CR^{d-1}$. Since
$\Delta(\chi\rho)=\chi\mu+2b\cdot\nabla\chi+\rho\Delta\chi$,
its total variation is at most $CR^{d-1}$. Integration by parts is
justified by uniform and local $H^1$ mollification of the continuous
function $v$. Bounding the right side by $C\delta R^{d-1}$ and splitting
the square over $A$ and $D\setminus A$ proves the estimate.
Finally $|b|\geq c$ almost everywhere, by Euler's identity and
$\rho(x)\geq c|x|$.
\end{proof}

\begin{corollary}[Exact rigidity]\label{cor:calibration-rigidity}
If $U_D\geq0$ and $U_D=0$ on $D^c$, then, up to null sets, $D=rK$ for
some $r\geq0$ and $U_D=2d\kappa(r-\rho)_+$.
\end{corollary}
\begin{proof}
If $|D|=0$, take $r=0$. Otherwise take $\delta=0$. Then $D=\{u>0\}$ up to null sets and
$\nabla u=-a\mathbf1_{\{u>0\}}\nabla\rho$.
Thus $u$ is Lipschitz and $u+a\rho$ is constant on each connected
component of its positivity set. On such a component write
$u=a(r-\rho)$. The segment from any point of the component toward the
origin cannot leave it: at a first exit its value would still be
strictly positive. Every component therefore contains the origin, so
there is only one. Its boundary lies in $\{\rho=r\}$; following each
radial segment now shows that the component is $rK$.
\end{proof}

\begin{corollary}[Few sites have small local time]\label{cor:low-occupation}
On the coarse event, let $\delta_n$ be a uniform bound for
$|\widetilde\ell_n-U_{D_n}|$. Then simultaneously for every $t\geq0$,
\[
 \#\{x\in A_n:\ell_n(x)\leq t\}
 \leq C(t+\delta_n)r_n^{d-1}.
\]
\end{corollary}
\begin{proof}
Apply the lemma with truncation threshold $t+\delta_n$ and $R=Cr_n$.
Every cell with $\ell_n(x)\leq t$ lies in $\{U_{D_n}\leq t+\delta_n\}$.
\end{proof}

\section{Limit laws}
\begin{proof}[Proof of Theorem~\ref{thm:main}\textup{(iv)}]
Put $b(x)=\kappa\lambda(x)\xi(x)$. Euler's identity gives $x\cdot b(x)=\kappa\lambda(x)\rho(x)$.
With $I_j$ indicating a first departure,
\[
 Q_n=|X_n|^2-n+2\kappa S_n,\qquad
 \Delta Q_j=2X_j\cdot(\Delta X_j+I_jb(X_j))
\]
is a martingale. Outside first departures its increment variance
is $4|X_j|^2/d$. Since $|b(x)|\leq1$,
\[
 \langle Q\rangle_n=\frac4d\sum_x\ell_n(x)|x|^2
             +O\!\left(\sum_{x\in A_n}|x|^2\right).
\]
The error is $O(r_n^{d+2})$. The cone and a Riemann sum give
\[
 \frac{\langle Q\rangle_n}{r_n^{d+3}}
 \longrightarrow8\kappa\int_{K}(1-\rho(z))|z|^2\,dz
       =\frac{8\kappa J_K}{d+3}\quad\hbox{almost surely}.
\]
The last identity is homogeneity. The bound
$|\Delta Q_j|\leq4|X_j|$ is predictable and verifies the
martingale Lindeberg and Stout increment conditions: its ratio to
$\sqrt{\langle Q\rangle_n/\log\log\langle Q\rangle_n}$ is
$O(\sqrt{\log\log n/n})$. Apply the martingale central limit theorem
\citep[Corollary~3.1]{HallHeyde1980} and iterated-logarithm law
\citep{Stout1970}.
Since $|X_n|^2=O(r_n^2)$ is negligible for $d\geq2$, divide by
$2\kappa$ to obtain the assertions for $S_n$.
These imply $S_n\sim n/(2\kappa)$ almost surely, and
differentiation gives the factor $1/(dV)$ for $\Rmom(n)$.
\end{proof}

\begin{theorem}[Local times centered by the cookie potential]
\label{thm:general-local-clt}
Under the assumptions of Theorem~\ref{thm:main}, use the Green
kernel $g$ of Lemma~\ref{lem:core-green}.
Define the range-dependent potential and normalization by
\[
 P_n(y)=\sum_{x\in A_n}\sum_{|e|=1}
       \left(\frac1{2d}-q_x(e)\right)g(x+e-y),\qquad
 s_n=\begin{cases}\sqrt{r_n\log r_n},&d=2,\\
                         \sqrt{r_n},&d\geq3.
             \end{cases}
\]
For every finite list $y_1,\ldots,y_k$,
$((\ell_n(y_i)-P_n(y_i))/s_n)_i$ tends to a centered Gaussian
vector of covariance
\[
 \Sigma_{ij}=\begin{cases}16\kappa/\pi,&d=2,\\
 2d\kappa(2G(y_i-y_j)-\mathbf1_{\{y_i=y_j\}}),&d\geq3.
 \end{cases}
\]
Write $\Sigma(y,y)$ for the corresponding diagonal variance.
At each fixed site, for either sign, almost surely,
\[
 \limsup_n\frac{\pm(\ell_n(y)-P_n(y))}
                   {s_n\sqrt{2\log\log n}}=\sqrt{\Sigma(y,y)}.
\]
Both assertions also hold with $U_{D_n}(y)$ in place of $P_n(y)$.
A deterministic cone requires the centering estimate below.
\end{theorem}
\begin{proof}
The Dynkin martingale $M^y$ for $g_y(x)=g(x-y)$ satisfies
\[
 \ell_n(y)-P_n(y)=-M_n^y+g(X_n-y)-g(-y).
\]
The endpoint is $O(\log n)$ in $d=2$ and $O(1)$ in $d\geq3$.
Define
\[
 \Gamma(f,h)(x)=\frac1{2d}\sum_{|e|=1}
 (f(x+e)-f(x))(h(x+e)-h(x))-\Delta f(x)\Delta h(x).
\]
The bound $|g(x+e)-g(x)|\leq C(1+|x|)^{1-d}$ shows that
first-cookie covariance corrections total $O_{y,z}(\log r_n)$
in the plane and $O_{y,z}(1)$ in higher dimensions. Thus
\[
 \langle M^y,M^z\rangle_n
       =\sum_x\ell_n(x)\Gamma(g_y,g_z)(x)+O_{y,z}(1+\log r_n).
\]
For $d\geq3$, summability, the cone profile and summation by parts
give $r_n^{-1}\langle M^y,M^z\rangle_n\to
2d\kappa(2G(y-z)-\mathbf1_{\{y=z\}})$.
In $d=2$, the covariance kernel is
$2/(\pi^2|x|^2)+O_{y,z}(|x|^{-3})$, whose sum over a ball of
radius $r$ is $(4/\pi)\log r+O(1)$.
The constant $4\kappa r_n$ of the cone gives
$(16\kappa/\pi)r_n\log r_n$; the term involving $\rho(x)$ and
the change from a Euclidean ball to a gauge ball cost only $O(r_n)$.
The profile error costs $o(r_n\log r_n)$. These are the bracket
limits. Bounded increments allow the multivariate martingale CLT
and the scalar LIL. The potential representation gives
$P_n(y)=U_{D_n}(y)+O(\log n)$ eventually almost surely,
which proves the assertion about continuum centering.
\end{proof}

\begin{lemma}[Centering at a fixed site with sparse weights]
Assume the Euclidean radius and volume estimates. Let $U_D$ denote the
\emph{unweighted} radial continuum potential, and set
\[
 q_n=\begin{cases}(\log n/r_n)^{1/2},&d=2,\\\log n/r_n,&d\ge3,
 \end{cases}\qquad
 w_n=\begin{cases}\sqrt{r_n}(\log n)^{5/2},&d=2,\\(\log n)^{d+1},&d\ge3.
 \end{cases}
\]
For every fixed $y$, almost surely for all large $n$,
\[
 U_{D_n}(y)-2d\kappa(r_n-|y|)-\frac{Q_n}{\omega_dr_n^d}
 =O_y(1+q_nw_n)
 =\begin{cases}O_y((\log n)^3),&d=2,\\O_y(1),&d\ge3.
 \end{cases}
\]
Consequently the actual local time satisfies
\[
 \ell_n(y)-2d\kappa(r_n-|y|)
 =-M_n^{g_y}+\frac{Q_n}{\omega_dr_n^d}
   +\begin{cases}O_y((\log n)^3),&d=2,\\O_y(\log n),&d\ge3.
   \end{cases}
\]
\end{lemma}
\begin{proof}
Write $\nu=1_{D_n}-1_{B(0,r_n)}$. Its support lies in
$\{v:||v|-r_n|\le Cw_n\}$ and $\|\nu\|_1\le Cr_n^dq_n$.
On this annulus,
\[
 \left||v|^{1-d}-r_n^{-d}|v|\right|
 \le Cr_n^{-d}||v|-r_n|.
\]
Integration against $\nu$ gives
\[
 U_{D_n}(0)-2d\kappa r_n
 =\frac{2\kappa}{\omega_dr_n^d}
   \left(\int_{D_n}|v|\,dv-\frac n{2\kappa}\right)+O(q_nw_n).
\]
Cell replacement costs $O(r_n^d)$. The sparse-weight condition gives
\[
 \left|S_n-\sum_{x\in A_n}|x|\right|
 \le Cr_n\sum_{|x|\le Cr_n}|\lambda_x-1|=O(r_n^d).
\]
Thus the parenthesis is $(Q_n-|X_n|^2)/(2\kappa)+O(r_n^d)$.
Since $|X_n|^2/r_n^d=O(1)$, this proves the assertion at the origin.
For fixed $y$ the difference between the kernels at $y$ and at zero
is $O_y(r_n^{-d})$ on the support of $\nu$, so its integral is $O_y(1)$.
The ball potential changes by exactly $-2d\kappa|y|$.
The local-time assertion follows from the Green-martingale
representation with its uniform $O(\log n)$ error.
\end{proof}

\paragraph{Deterministic centering in the Euclidean case.}
The preceding lemma completes Theorem~\ref{thm:radial}\textup{(ii)}
from Theorem~\ref{thm:general-local-clt}. Indeed, the moment LIL gives
\[
 Q_n/r_n^d=O\bigl(r_n^{(3-d)/2}\sqrt{\log\log n}\bigr)
 \quad\hbox{almost surely},
\]
which, together with the centering error, is negligible on both the
fixed-site CLT and LIL scales. The same Green martingales therefore
give the asserted joint Gaussian law and signed iterated-logarithm laws.


\section{Sharpness and fluctuation lower bounds}

We continue with $\rho(x)=|x|$ and put $u_x=x/|x|$ for
$x\ne0$, $u_0=0$.
Write $r=r_n$, $L=\log n$, $R_-=\Rin(n)$, and $R_+=\Rout(n)$.
All estimates below allow arbitrary first-cookie covariances.
We use the already proved quantitative shape estimates on an event
$\Omega_n$, for a fixed constant $K_0<\infty$, with $\mathbb P(\Omega_n^c)\leq Cn^{-10}$:
\begin{align}
 R_+&\leq K_0r,&
 \sup_x|\ell_n(x)-2d\kappa(r-|x|)_+|
      &\leq Cr^{1/2}L^{3/2},\label{eq:lower-input-profile}\\
 |R_--r|+r^{1-d}|D_n\mathbin\triangle B_r|
      &\leq C h_d,&
 h_d&=\begin{cases}\sqrt{rL},&d=2,\\L,&d\geq3.
                  \end{cases}\label{eq:lower-input-volume}
\end{align}
The discrete-to-continuum potential estimate and Dynkin's identity
then give, uniformly for $|y|\leq r/2$,
\begin{equation}\label{eq:lower-bulk-representation}
 \ell_n(y)-2d\kappa(r-|y|)=-M_n^y+O(h_d+L)
 \quad\hbox{on }\Omega_n.
\end{equation}
Indeed, $E=D_n\setminus B_{R_-}$ lies a distance of order $r$
from such $y$, so its drift potential has absolute value at most
$Cr^{1-d}|E|\leq Ch_d$; the potential of $B_{R_-}$ is the cone
of height $2d\kappa R_-$, and the remaining error is $O(L)$.
\subsection{One exponential-martingale lemma}

\begin{lemma}\label{lem:lower-exp}
Fix $0<c_0\leq C_0$.  Let $S^1,\ldots,S^m$ be real
martingales starting at zero with increments bounded by $b$.
Suppose an event $\Omega$, of probability at least $1-\alpha$,
satisfies $c_0\leq\langle S^i\rangle_n\leq C_0$ for every $i$.
There are $c,C>0$, depending only on $c_0,C_0$, such that if
$\beta\geq C$ and $\beta b\leq c$, then
\[
 \mathbb P(\pm S_n^i\geq c\beta)
       \geq\tfrac14e^{-C\beta^2}-\alpha.
\]
If additionally
$|\langle S^i,S^j\rangle_n|\leq\delta$ on $\Omega$ for $i\ne j$
and $\beta^2\delta+\beta^3b\leq1$, then
\[
 \mathbb P\{\max_i S_n^i<c\beta\}
 \leq C\left(m^{-1}e^{C\beta^2}+\beta^2\delta+
                    \beta^3b+e^{C\beta^2}\sqrt\alpha\right).
\]
\end{lemma}
\begin{proof}
Stop all martingales just before any bracket would exceed $2C_0$;
this changes none of them on $\Omega$.  For a martingale $S$ whose
increments are at most $b_1$ and whose bracket is bounded by $8C_0$,
put
\[
 K_S(s)=\sum_{j<n}\log\mathbb E(e^{s\Delta S_j}\mid\mathcal F_j),
 \qquad Z_S(s)=e^{sS_n-K_S(s)}.
\]
Conditional Taylor expansion, for $|s|b_1\leq1/2$, gives
\[
 K_S(s)=\tfrac12s^2\langle S\rangle_n+O(|s|^3b_1),
 \quad \mathbb E Z_S(s)=1,
 \quad \mathbb E Z_S(s)^2\leq e^{Cs^2}.
\]
For the last bound use
$Z_S(s)^2=Z_S(2s)e^{K_S(2s)-2K_S(s)}$.
Apply these formulas to $S^i$ and $S^i+S^j$, with $s=\beta$,
and abbreviate $Z_i=Z_{S^i}(\beta)$.
On $\Omega$, $Z_i\geq1/2$ implies $S_n^i\geq c\beta$,
since $K_{S^i}(\beta)\geq c_0\beta^2/4$.
Paley--Zygmund therefore gives the first assertion; use $-S^i$
for the opposite sign.
For $i\ne j$,
\[
 Z_iZ_j=Z_{S^i+S^j}(\beta)
       \exp\{\beta^2\langle S^i,S^j\rangle_n+O(\beta^3b)\}.
\]
On $\Omega$ the exponential is at most
$1+C(\beta^2\delta+\beta^3b)$; on its complement apply
Cauchy--Schwarz and the second-moment bound.  Consequently
\[
 \mathbb E Z_iZ_j\leq1+C(\beta^2\delta+\beta^3b)
                              +e^{C\beta^2}\sqrt\alpha.
\]
The average $\overline Z=m^{-1}\sum_i Z_i$ has mean one and
variance bounded by the right side of the asserted probability
estimate.  If $\overline Z\geq1/2$ on $\Omega$, some
$S_n^i\geq c\beta$.  Chebyshev, with the additional error $\alpha$,
finishes the proof.
\end{proof}

\subsection{Planar radii and width}

Let $d=2$.  Define the weighted moment and vector martingales
\[
 S_n=\sum_{x\in A_n}\lambda_x|x|,\qquad
 Q_n=|X_n|^2-n+2\kappa S_n,\qquad
 Z_n=X_n+\kappa\sum_{x\in A_n}\lambda_xu_x.
\]
On $\{R_+\leq K_0r\}$ the defect assumption gives
\begin{equation}\label{eq:lower-defect-scores}
 \left|S_n-\sum_{x\in A_n}|x|\right|\leq Cr^2,
 \qquad
 \left|\sum_{x\in A_n}(\lambda_x-1)u_x\right|\leq Cr.
\end{equation}
The first-cookie covariance corrections contribute $O(r^4)$ to
$\langle Q\rangle_n$ and $O(r^2)$ to each bracket of $Z$.
Thus \eqref{eq:lower-input-profile} gives, on $\Omega_n$,
\begin{equation}\label{eq:lower-QZ-brackets}
 \langle Q\rangle_n=(4\pi\kappa/5+o(1))r^5,
 \qquad
 \langle v\cdot Z\rangle_n=|v|^2n/2+O(|v|^2r^2).
\end{equation}
These formulas also hold almost surely with relative error $o(1)$
by the qualitative shape theorem.

The unweighted moment radius
\[
 T_n=\left(\frac3{2\pi}\sum_{x\in A_n}|x|\right)^{1/3}
\]
satisfies $R_-\leq T_n+C$ and $R_+\geq T_n-C$ by comparison
with unit cells.  The weighted moment radius differs from $T_n$
by $O(1)$ on $\Omega_n$ and eventually almost surely, by
\eqref{eq:lower-defect-scores} and the shape theorem.
The general moment LIL therefore proves
\begin{equation}\label{eq:lower-radii-lil}
 \limsup_n\frac{r-R_-}{\sqrt{r\log\log n}},\quad
 \limsup_n\frac{R_+-r}{\sqrt{r\log\log n}}
 \ \geq\frac1{\sqrt{10\pi\kappa}}\quad\hbox{almost surely}.
\end{equation}
To obtain polynomial lower tails, stop $Q$ on exiting $B_{K_0r}$
and apply Lemma~\ref{lem:lower-exp} to $Q/r^{5/2}$.
Its increments are bounded by $4K_0r^{-3/2}$, and its bracket lies
between positive constants on $\Omega_n$ by
\eqref{eq:lower-QZ-brackets}.
Choose $\beta=a_p\sqrt{\log n}$, with $a_p>0$ small enough that
$Ca_p^2<\min(p,10)/2$.  Each sign then has probability exceeding
$n^{-p}$ after intersecting with $\Omega_n$ and subtracting its
failure probability once more, for large $n$.
Since $|X_n|^2$ and the defect in \eqref{eq:lower-defect-scores}
are $O(r^2)=o(r^{5/2}\sqrt{\log n})$, the identity for $Q$ and
factoring $T_n^3-r^3$ yield, with a suitable $c_p>0$,
\begin{equation}\label{eq:lower-radii-tails}
 \mathbb P\{r-R_-\geq c_p\sqrt{r\log n}\}\geq n^{-p},
 \qquad
 \mathbb P\{R_+-r\geq c_p\sqrt{r\log n}\}\geq n^{-p}.
\end{equation}

For the width set $V_n=\sum_{x\in A_n}u_x$.
The sum over the lattice disk $\{|x|<R_-\}$ vanishes by symmetry.
When $R_-,R_+\asymp r$, integration of the positive part of any
directional component over the remaining annulus gives
\begin{equation}\label{eq:lower-width-geometry}
 |V_n|\leq R_+^2-R_-^2+Cr,
 \qquad |Z_n|\leq\kappa(R_+^2-R_-^2)+Cr.
\end{equation}
Indeed, $\int_0^{2\pi}(\cos\theta)_+\,d\theta=2$; replacing
each unit vector by its values in the corresponding cell costs
$O(r)$, as does the weighted correction.
The martingale $e_1\cdot Z$ has increments at most two.
Apply Lemma~\ref{lem:lower-exp} to $(e_1\cdot Z)/\sqrt n$,
using its bracket bounds on $\Omega_n$ from
\eqref{eq:lower-QZ-brackets}.  Equation
\eqref{eq:lower-width-geometry} and $n=(4\pi\kappa/3)r^3$ give
\begin{equation}\label{eq:lower-width-tail}
 \mathbb P\{R_+-R_-\geq c_p\sqrt{r\log n}\}\geq n^{-p}.
\end{equation}
The same bracket limit, the bounded-increment martingale LIL,
and $R_-+R_+\sim2r$ give
\begin{equation}\label{eq:lower-width-lil}
 \limsup_n\frac{R_+-R_-}{\sqrt{r\log\log n}}
       \geq\sqrt{\frac\pi{3\kappa}}\quad\hbox{almost surely}.
\end{equation}
Finally the vector martingale CLT yields
$Z_n/\sqrt n\Longrightarrow N(0,I_2/2)$.
On $\{R_+-R_-\leq a\sqrt r\}$,
\eqref{eq:lower-width-geometry} gives
$|Z_n|/\sqrt n\leq\sqrt{3\kappa/\pi}\,a+o(1)$ almost surely.
Since $\mathbb P\{|N(0,I_2/2)|\leq t\}=1-e^{-t^2}$,
\begin{equation}\label{eq:lower-width-small}
 \limsup_n\mathbb P\{R_+-R_-\leq a\sqrt r\}
       \leq1-e^{-3\kappa a^2/\pi}\qquad(a\geq0).
\end{equation}

\subsection{The maximum of the bulk local times}

Let $m=\lfloor r^{1/8}\rfloor$,
$y_i=i\lceil r^{1/4}\rceil e_1$, and put
\[
 f_i=\begin{cases}g_{y_i+ke_1}-g_{y_i},\quad k=\lfloor r^{1/8}\rfloor,
                         &d=2,\\g_{y_i},&d\geq3,
       \end{cases}
 \qquad
 s=\begin{cases}\sqrt{r\log r},&d=2,\\\sqrt r,&d\geq3,
      \end{cases}
 \qquad S^i=-M^{f_i}/s.
\]
All the poles lie in $B_{\sqrt r}$ for large $n$.
We verify, uniformly on $\Omega_n$,
\begin{equation}\label{eq:lower-separated}
 c_0\leq\langle S^i\rangle_n\leq C_0,
 \qquad |\langle S^i,S^j\rangle_n|\leq Cr^{-1/4}\quad(i\ne j).
\end{equation}
This calculation uses only the simple-random-walk covariance kernel
$\Gamma$ and the cone estimate; no independence is claimed.
Write $W_{ij}=\sum_x|\Gamma(f_i,f_j)(x)|$ and
$T_{ij}=\sum_x\min(|x|,r)|\Gamma(f_i,f_j)(x)|$.
The Green-kernel gradient bound and $2ab\leq a^2+b^2$ give
\[
 (W_{ij},T_{ij})\leq
 \begin{cases}(C\log r,C\sqrt r\log r),&d=2,\\
                  (C,C(\sqrt r+\log r)),&d\geq3.
       \end{cases}
\]
For completeness, in $d\geq3$ the squared increments are bounded
by $C(1+|x-y_i|)^{2-2d}$; its sum is finite and its first truncated
moment is $O(\log r)$ in $d=3$ and $O(1)$ in larger dimensions.
In $d=2$, an increment of $f_i$ is bounded by the sum of reciprocal
distances to its two poles within distance $3k$, and by
$Ck(1+|x-y_i|)^{-2}$ outside.  Squaring and summing gives
$O(\log r)$ and a first moment $O(k\log r)$, respectively.
Translation of the poles costs at most $C\sqrt r$ times the
unweighted sums.  The finitely supported Laplacian terms in
$\Gamma$ satisfy the same bounds.

Replacing first-cookie covariances by $\Gamma$ costs $O(1)$ in
$d\geq3$ and $O(\log r)$ in $d=2$.
Replacing $\ell_n$ by its cone and then by the constant $2d\kappa r$
costs at most $Cr^{1/2}L^{3/2}W_{ij}+2d\kappa T_{ij}$.
After division by $s^2$ all these errors are
$O(r^{-1/2}L^{3/2})=o(r^{-1/4})$.
Thus the bracket differs by this error from
\[
 \Sigma_{ij}=\frac{2d\kappa r}{s^2}
                         \sum_x\Gamma(f_i,f_j)(x).
\]
For $d\geq3$, summation by parts gives
$\Sigma_{ij}=2d\kappa(2G(y_i-y_j)-\mathbf1_{\{i=j\}})$.
Its diagonal is positive, and its off-diagonal entries are
$O(r^{(2-d)/4})$.
For $d=2$, writing $z=y_i-y_j$, the same identity gives
\[
 \sum_x\Gamma(f_i,f_j)(x)
   =2\{g(z+ke_1)+g(z-ke_1)-2g(z)\}-2\mathbf1_{\{i=j\}}.
\]
The diagonal of $\Sigma$ tends to $4\kappa/\pi$, because
$g(ke_1)=(2/\pi)\log k+O(1)$.
For $i\ne j$, the displayed second difference is
$O(k^2/|z|^2)$, giving an off-diagonal bound $O(r^{-1/4})$.
This proves \eqref{eq:lower-separated}.

The normalized martingale increments are bounded by $Cr^{-1/2}$.
Apply Lemma~\ref{lem:lower-exp} with
$\beta=a\sqrt{\log r}$, where $a>0$ is so small that
$Ca^2<1/16$.  Since $m\asymp r^{1/8}$,
$\delta=Cr^{-1/4}$, and $\alpha=O(n^{-10})$, the probability
that all $S_n^i<c\sqrt{\log r}$ is at most $Cr^{-1/16}$.
The error in \eqref{eq:lower-bulk-representation} is
$o(s\sqrt{\log r})$.  In $d\geq3$ a large $S_n^i$ gives a large
positive local-time deviation at $y_i$; in $d=2$ it gives a large
difference of the deviations at its two poles, so one has large
absolute value.  Consequently, for some $c>0$,
\begin{equation}\label{eq:lower-bulk-max}
 \mathbb P\left\{\max_{|y|\leq\sqrt r}
 |\ell_n(y)-2d\kappa(r-|y|)|\geq
 c\begin{cases}\sqrt r\log n,&d=2,\\\sqrt{r\log n},&d\geq3
   \end{cases}\right\}\geq1-n^{-c}
\end{equation}
for all sufficiently large $n$.

\subsection{Logarithmic radii under ellipticity}

For this paragraph assume additionally $q_x(e)\geq p_0>0$
for every $x,e$; later departures already have probability $1/(2d)$.
There are $c,C,h_0>0$ such that, for $h_0\leq h\leq r/8$,
\begin{align}
 \mathbb P\{R_-\geq r-h,\ R_+\leq r+h\}
   &\leq e^{-cr^{d-1}e^{-Ch}},\label{eq:lower-log-joint}\\
 \mathbb P\{R_+\leq r+h\}
   &\leq e^{-cr^{d-1}e^{-Ch}}+e^{-cr^{d-3}h^2}.
       \label{eq:lower-log-outer}
\end{align}
To prove these estimates, from any $x$ with $|x|>r-2h$ take
$T=\lceil4h\sqrt d\rceil$ steps in the signed coordinate of
largest $|x_i|$.  The Euclidean distance increases at speed at least
$1/\sqrt d$, so this path exits $B_{r+h}$; its conditional
probability is at least $p_0^T$ after decreasing $p_0$ if necessary.
Among $k$ distinct arrivals at sites of any fixed
$\Lambda\subset\{|x|>r-2h\}$ by time $n$, at least
$\lfloor(k-T)/(T+1)\rfloor$ disjoint such trials can be started
and completed by time $n$.  Successive conditioning bounds their
joint failure probability by
$(1-p_0^T)^{\lfloor(k-T)/(T+1)\rfloor}$.
If $R_-\geq r-h$, the annulus
$r-2h<|x|<r-h-1$ supplies $k\geq cr^{d-1}h$ sites, proving
\eqref{eq:lower-log-joint}.

For \eqref{eq:lower-log-outer}, let
$\Lambda=\{r-2h<|x|\leq r+h\}$.
The same trials apply if $|A_n\cap\Lambda|\geq c_1r^{d-1}h$.
Otherwise, on $R_+\leq r+h$, cell comparison and the defect bound
give
\[
 \sum_{x\in A_n}\lambda_x|x|
 \leq\frac{n}{2\kappa}-c_2r^dh
                     +2r|A_n\cap\Lambda|+Cr^d.
\]
Choose $c_1$ small and then $h_0$ large.  The identity
$Q_n=|X_n|^2-n+2\kappa\sum\lambda_x|x|$
forces $Q_n\leq-cr^dh$; the endpoint $|X_n|^2$ is absorbed
by $h_0$ in $d=2$ and by large $r$ in higher dimensions.
Stop $Q$ on leaving $B_{r+h}$.  Its increments are bounded by
$4(r+h)$, so Azuma's inequality gives
$\mathbb P\{Q_n\leq-cr^dh,\ R_+\leq r+h\}
\leq e^{-cr^{d-3}h^2}$, proving the second estimate.
Taking $h=c\log n$ with $c>0$ sufficiently small and applying
Borel--Cantelli gives, almost surely for all sufficiently large $n$,
\[
 \max(r-R_-,R_+-r)>c\log n;
 \qquad R_+-r>c\log n\quad\hbox{if }d\geq3.
\]
The ellipticity assumption is used only in the prescribed-path
trials and is not needed for any preceding lower bound.


\paragraph{Draft and verification scope.}
This is a new mathematical draft. The Lean certificate accompanying
\cite{BouRabeePeres2026CERW} is pinned to that paper and does not certify
the generalized statements here.
\bibliographystyle{plainnat}
\bibliography{references}
\end{document}
